\section{A concrete zero-credit theorem for ZC215}\label{sec:credit}

\subsection{Why this is a distinct simplification}

ZC215 uses DROP tag 0, NIP tag 1, MP tag 2, and AX tags at least 3. Its
proposed refill enumerates every raw candidate in order $0,1,2,\ldots$.
The ordinary proof in the repository gives an injective map from NIP
occurrences to DROPs in strictly earlier candidates \cite{zerocredit}.
The existing GC--ZC abstract prototype assumes the key lowering and
injectivity properties. The new \code{CantorCredit.v} proves those
properties for the actual codec, proves total unique decoding, and constructs
the literal schedule prefix. It therefore closes a specific formal gap;
it does not claim discovery of the injection itself.

In the project's command semantics, DROP pushes one theorem head and NIP
removes one. The conservative credit assignment is
\[
 \mathrm{DROP}=+1,\qquad \mathrm{NIP}=-1,
 \qquad\mathrm{MP}=\mathrm{AX}=0.
\]
The current AX can add further heads, but the proof below does not need them.
That is what makes it useful for future source simplification.

\subsection{Exact codec, including zero}

Put
\[
 P(x,y)=\frac{(x+y)(x+y+1)}2+y.
\]
The formal file uses the standard-library Cantor pairing with this
orientation. Complete command fields are:
\begin{center}
\begin{tabular}{@{}ll@{}}\toprule
Command & Field sequence\\\midrule
DROP & $[0]$\\
NIP & $[1]$\\
MP$(x)$ & $[2,x]$\\
AX$(c;x,y)$ & $[c+3,y,x]$\\\bottomrule
\end{tabular}
\end{center}
For fields $a_1,\ldots,a_r$, define
\[
 E([a_1,\ldots,a_r],z)=P(a_1,P(a_2,\ldots,P(a_r,z)\ldots)).
\]
For a nonempty command word $L$, its raw code is defined backward by
\begin{align*}
 \rho([k])&=E(\operatorname{fields}(k),0),\\
 \rho(k::L)&=E(\operatorname{fields}(k),1+\rho(L))
 \quad(L\ne[]).
\end{align*}
In particular $\rho([\mathrm{DROP}])=0$. Zero is a valid nonempty raw
candidate; a terminal DROP in a longer word is distinguished by the
positive continuation at its predecessor. Discarding this zero alias would
break the exact scheduler argument.

\begin{lemma}[Concrete decoding]
Every $n\in\N$ is the code of a unique nonempty command word $W(n)$.
The supplied executable decoder satisfies
\[
 \rho(W(n))=n,\qquad W(\rho(L))=L.
\]
\end{lemma}
\begin{proof}
Invert Cantor pairing for the outer tag, then consume precisely the fields
prescribed by that tag. This yields a unique command and residual $r$.
For $r=0$, the word terminates. For $r>0$, recurse on $r-1$.
The residual is at most the enclosing code, so $r-1<n$ in the recursive
case. Strong induction proves existence; command-field injectivity and
the successor continuation convention prove uniqueness. The concrete
fuel-based decoder uses $n+1$ as a sufficient recursion bound. Its default
fuel-exhaustion branch is proved unreachable under that bound.
\end{proof}

No formula-decoding or theoremhood assumption occurs here. The result
concerns all raw natural codes, exactly as needed when malformed formula
codes or numerical aliases are permitted by the numerical calculus.

\subsection{Occurrence injection and prefix safety}

\begin{theorem}[Strict lowering of a selected NIP]
Let command position $j$ of a nonempty word $L$ be NIP. Let $L'$ be
obtained by replacing just that command with DROP. Then
\[
 \rho(L')<\rho(L),
\]
and position $j$ of $W(\rho(L'))$ is DROP.
\end{theorem}
\begin{proof}
The replacement preserves the command arity, all other fields, and every other
command. At the changed tag, with residual $r$, the exact decrease is
\[
 P(1,r)-P(0,r)=r+1>0.
\]
Every surrounding successor and pairing in the earlier command prefix
preserves strict inequality. The round-trip decoder theorem supplies the
claim about the same position in the lowered word.
\end{proof}

Retaining the position is essential: the objects counted are occurrences,
not merely command values or candidate indices. Map a NIP occurrence
$(L,j)$ to $(L',j)$. This map is injective, because replacing the image
DROP at that same position by NIP recovers $L$ exactly.

\begin{theorem}[Concrete exhaustive prefix credit]\label{thm:credit}
For every $n,j\in\N$, consider the schedule prefix consisting of
\[
 W(0)W(1)\cdots W(n-1)
 \quad\hbox{followed by the first $j$ positions of }W(n).
\]
Its number of NIP occurrences is at most its number of DROP occurrences.
If the next actual command is NIP, the strict inequality already holds
immediately before that command.
\end{theorem}
\begin{proof}
Each included NIP is mapped injectively to a DROP in a strictly smaller raw
candidate. That candidate is complete in the same prefix, so all image
occurrences are already included. Injectivity proves the weak inequality.
Apply it to the prefix including one further NIP: its NIP count increases
by one, while its DROP count does not. This yields the strict inequality
for the prefix immediately before that NIP.
\end{proof}

The formal weak inequality is \code{exhaustive\_schedule\_zero\_credit}.
The strict version is \code{exhaustive\_credit\_before\_nip}.
They have no abstract scheduling premises.
The concrete finite list \code{schedule\_cut n j} contains all
positions of the earlier candidates and $j$ positions of the current one.
If $j$ exceeds its length, extra positions have no command and contribute
neither kind; ordinary execution prefixes use the actual length bound.

\subsection{Application and remaining machine work}

Represent the current numerical stack as a list of explicit theorem heads
above a fixed arbitrary opaque natural tail $J$. Begin with no explicit
heads. DROP increases the protected length, NIP decreases it, and the
remaining commands conservatively do not decrease it. The strict part of
Theorem~\ref{thm:credit} ensures a protected head exists whenever a NIP is
about to execute. Thus it never inspects $J$. In the actual current ZC
proposal, this includes the initialized tail $J=5$ without asserting that
5 itself encodes a theorem stack \cite{zerocredit}.

This argument applies to executed commands before the first halt once the
machine/source bridge establishes the literal schedule: all earlier
candidates have then completed without halting. It cannot justify a physical
implementation that skips a candidate or returns to the wrong control.
The new file does not model that physical implementation. Its remaining
interfaces are exact parsing by the source, exact refill order, command
soundness and strong completeness, positive binary simulation, startup,
and the inherited numerical-calculus/ZFC calibration.

\code{CantorCredit.v} has 408 lines. It compiled, printed nine closed
assumption reports, and passed recursive \code{coqchk}. These results
make the ZC scheduler argument a reusable concrete theorem. They do not
remove the obligations needed for a complete ZC215 endpoint, and they do
not reduce the number of physical states. In particular, the present
primary K215's zero initialized stack already avoids this arbitrary-tail
credit argument; adopting ZC solely to shorten that existing endpoint
would not automatically be beneficial.
